\begin{claim}{B.9}{Proved}{Uniformly finite stages give a constant finite category}
Let $m\geq0$ be an integer and let $K:I\to\Cathot_\U$ be
a $\U$-small filtered diagram with $|\Mor K_i|\leq m$ for
every $i$. There is a category $L$ with at most $m$ arrows and
an isomorphism
\[
                 (K_i)_{i\in I}\cong L
                   \quad\hbox{in }\Ind_\U(\Cathot_\U),
\]
where the right side is the constant ind-object. In particular,
the bounded-test invariance of such an object is the invariance
of this finite representable. On finite test categories its
membership is given by the already proved criterion B.8.
\end{claim}
\begin{proof}
The case $m=0$ has only the empty category and is immediate.
Choose the full subcategory $\mathcal D$ of $\Cathot_\U$
with one representative of each isomorphism type of category
having at most $m$ arrows. There are finitely many such finite
structures: both the object and arrow sets have bounded finite
cardinality, and their source, target, identity and composition
operations have finitely many possibilities. Every functor set,
and hence every categorical homotopy-class set, is finite.
Thus $\mathcal D$ has finitely many objects and arrows.
Transporting each stage across a chosen isomorphism makes $K$
a diagram in $\mathcal D$.

Put $F=\operatorname{colim}_i h_{K_i}$ on $\mathcal D$.
For each $A\in\mathcal D$ its value has cardinality at most
\[
             M_A=\max_{T\in\mathcal D}|\mathcal D(A,T)|.
\]
Indeed, any $M_A+1$ distinct elements would have representatives
at one common stage by filteredness. Those representatives
would have to be distinct there, contradicting the definition
of $M_A$. The category of elements $\mathcal E$ of $F$ is
therefore finite. It is filtered: its objects and arrows are
represented at stages of the original filtered diagram, so
finitely many representatives have a common stage and equality
of their colimit images is witnessed at a later stage.

A filtered category has a cocone on each finite diagram. Apply
this to the identity diagram of the finite category $\mathcal E$.
It gives an object $e=(T,x)$ and arrows
$t_d:d\to e$ with $t_e t_d=t_d$ for every $d$ and in particular
$t_e^2=t_e$. The density presentation
$F=\operatorname{colim}_{d=(A,a)\in\mathcal E}h_A$ gives maps
\[
             h_T\xrightarrow{\rho}F
                       \xrightarrow{\sigma}h_T,
             \qquad \rho\sigma=1_F,
             \quad\sigma\rho=h_{t_e}.
\]
Here $\rho$ is its canonical map at $e$ and $\sigma$ is induced
by the cocone $t_d$; the displayed identities follow on each
representable summand from the cocone equations. The full
inclusion $\mathcal D\hookrightarrow\Cathot_\U$ induces a
fully faithful inclusion on ind-categories, as is immediate from
their limit--colimit Hom formula. These maps therefore exhibit
the original ind-object as the image of an idempotent in
$\Hom_{\Cathot}(T,T)$.

Choose an actual endofunctor $f:T\to T$ representing that
idempotent class. The set of endofunctors of $T$ is finite, so
the sequence of powers of $f$ is eventually periodic. Choose a
positive integer $r$ beyond the preperiod and divisible by its
period. Then $f^{2r}=f^r$ strictly, while $[f^r]=[f]$ because
$[f]^2=[f]$. Let $L$ have the objects and arrows fixed by $f^r$.
They are closed under identities and composition. Its inclusion
$j:L\to T$ and the functor $s:T\to L$ given by $f^r$ satisfy
\[
                         sj=1_L,\qquad js=f^r.
\]
This splits the original idempotent in $\Cathot$, with
$|\Mor L|\leq|\Mor T|\leq m$. Combining the two splittings
gives the required ind-isomorphism. Evaluation on every test
category preserves that isomorphism; the final invariance
assertion follows by this evaluation and B.8.
\end{proof}
